\documentclass[10pt,a4paper]{amsart}
\usepackage[margin=19mm]{geometry}
\usepackage{amsmath,amssymb,mathtools}
\usepackage[scr=rsfs,cal=euler]{mathalfa}
\usepackage{xfrac}
\usepackage[unicode,hidelinks,bookmarks=false]{hyperref}
\newcommand{\ie}{i.e.\ }
\newcommand{\cf}{cf.\ }
\newcommand{\resp}{respectively\ }
\newcommand{\Subsection}{\subsection}
\providecommand{\nobreakdash}{\nobreak\hbox{-}}
\setlength{\emergencystretch}{2em}
\multlinegap0pt
\usepackage{bm}
\usepackage{bbm}
\usepackage{dsfont}
\usepackage{xspace}
\usepackage{cancel}
\usepackage{mathtools}
\usepackage{amsthm}
\makeatletter
\@ifundefined{lemma}{\newtheorem{lemma}{Lemma}}{}
\@ifundefined{proposition}{\newtheorem{proposition}{Proposition}}{}
\makeatother
\title[The smallest singular value of sparse discrete random matric]{The smallest singular value of sparse discrete random matrices}
\author{Kexin Yu}
\date{Source: arXiv:2507.17741v1 (CC BY 4.0)}
\begin{document}
\maketitle

\noindent\emph{Source and licence.} The smallest singular value of sparse discrete random matrices. Version arXiv:2507.17741v1 (CC BY 4.0), distributed under the Creative Commons Attribution 4.0 International licence, as recorded in the arXiv licence metadata for this exact version and shown on the abstract page https://arxiv.org/abs/2507.17741v1. Full rights notice: licenses/arxiv-2507.17741-CC-BY-4.0.txt.\par\smallskip

\noindent\textit{Editorial scope.} Counted unit: the Proposition whose source label is \texttt{pro1} in the article (source0.tex lines 451--453, source label \texttt{pro1}), delivered together with the article's own complete original proof of it (source0.tex lines 459--501, one \texttt{proof} environment of 43 lines). No mathematics is written, repaired or re-derived: statement and proof are the article's own text line for line. The proof is detached from the statement in the source: the article states the result at lines 451--453 and prints its proof at lines 459--501, and the proof environment's own header names the statement, so the association is the article's own. Required context retained uncounted: lcd (lines 239--252). Every definition, display and label of those lines is retained unchanged. Omitted from this package and not redistributed: 1--238, 253--450, 454--458, 502--713. Nothing inside a retained excerpt is omitted or altered: every retained body is delivered exactly as the article's lines are written, so no omission receipt of any kind accompanies this package. Citations: the retained excerpts carry no citation command at all, so no bibliography entry of the article is transcribed and no reference is redistributed. Reference adaptation: none was needed and none was made; every reference inside the delivered unit resolves inside it, so no unresolved reference or citation is produced. Printed slips kept literal: no printed word, symbol, hypothesis or number of the counted statement or of its proof was altered. Layout and class: the article's own class is \texttt{article} with packages including psfrag, parskip, float, graphicx, epsfig, amsmath, amsthm, amssymb, verbatim, multicol, url, latexsym, mathrsfs, hyperref; the wrapper is the lane's pinned \texttt{amsart} editorial wrapper at 10pt on A4, which restates the theorem-like environments the retained text needs and adds the article's own macro definitions that the retained text needs (none), with extra wrapper packages (bm, bbm, dsfont, xspace, cancel, mathtools). The delivered numbering is the unit's own. Assets: no retained span contains a figure environment, a graphics-inclusion command, a PDF-page inclusion, a table or a TikZ picture; no image file is copied into this package, the package declares no asset dependency and no image file is redistributed with it. Licence: version arXiv:2507.17741v1 of this article is distributed under the Creative Commons Attribution 4.0 International licence, as recorded in the arXiv licence metadata for this exact version and shown on the abstract page https://arxiv.org/abs/2507.17741v1. A full rights notice is delivered at licenses/arxiv-2507.17741-CC-BY-4.0.txt. Characters: every retained excerpt is pure ASCII, so the delivered engine, XeLaTeX with its default Latin Modern text font, reports no missing character.
\begin{lemma}\label{lcd}
Let $\xi_1, \ldots, \xi_n$ denote i.i.d. $\mu-$lazy random variables, $0<\mu_{n}\leq1/2$. Consider a unit vector $\boldsymbol{a}=\left(a_1, \ldots, a_n\right) \in \mathbb{S}^{n-1}$. Let
\begin{align*}
S:=\sum_{i=1}^n \xi_i a_i.
\end{align*}
Then, for every $\alpha>0$, and for
\begin{align*}
\delta \geq \frac{(2 / \pi)}{\operatorname{LCD}_{\gamma, \alpha}(\boldsymbol{a})},
\end{align*}
we have
\begin{align*}
\mathcal{L}(S, \delta) \lesssim \frac{\delta}{\sqrt{\mu_{n}}\gamma}+\exp \left(-2\mu_{n}\alpha^2\right).
\end{align*}
\end{lemma}

\begin{proposition}\label{pro1}
$\textbf{P}\left(\exists \boldsymbol{a} \in \Gamma^1(\eta):\left\|M_n \boldsymbol{a}\right\|_2 \leq \eta\right) \lesssim \frac{\eta n^{3 / 2}}{\sqrt{\mu_{n}}}+n \exp (-2\sqrt{n}\mu_{n})$.
\end{proposition}

\begin{proof}[Proof of Proposition \ref{pro1}]
For every $M_n$ satisfying the above assumptions, since $M_n$ and $M_n^T$ have the same singular values, the event $s_n(M_n)\leq\eta$ implies the existence of a unit vector $\boldsymbol a'\in\mathbb R^n$ for which
\begin{align*}
\|\boldsymbol a'{}^T M_n\|_2 \leq\eta.
\end{align*}
For each fixed $M_n$, we define the unit vector $\boldsymbol a'\in\mathbb R^n$ satisfying the above condition to be $\boldsymbol a'_{M_n}$. Since the matrix rows are independent of each other, this leads to a partition of the space of all the above types of matrices with the smallest singular value at most $\eta$. Then, by taking a union bound, 
\begin{align*}
\textbf{P}\left(\exists \boldsymbol{a} \in \Gamma^1(\eta):\left\|M_n \boldsymbol{a}\right\|_2 \leq \eta\right)\leq n \textbf{P}\left(\exists\,\boldsymbol a\in\Gamma^1(\eta):
\|M_n\boldsymbol a\|_2\le\eta \wedge
\|\boldsymbol a'_{M_n}\|_\infty=|a'_n|\right),
\end{align*}
it suffices to show the following:
\begin{align}\label{pro30}
\textbf{P}\left(\exists\,\boldsymbol a\in\Gamma^1(\eta):
\|M_n\boldsymbol a\|_2\le\eta \wedge
\|\boldsymbol a'_{M_n}\|_\infty=|a'_n|\right)
\lesssim\eta\sqrt {\frac{n}{\mu_{n}}} +e^{-2\sqrt n\mu_{n}}.
\end{align}
To establish this, reveal the first \(n-1\) rows \(X_1,\dots,X_{n-1}\).  If there are some $\boldsymbol a\in\Gamma^1(\eta)$ satisfies $\|M_n\boldsymbol a\|_2\le\eta$, then based solely on \(X_1,\dots,X_{n-1}\) one can select a vector \(\boldsymbol y\in\Gamma^1(\eta)\) such that
\[
\left(\sum_{i=1}^{n-1}(X_i\cdot\boldsymbol y)^2\right)^{1/2}
\le\eta.
\]
Thus, for any vector $\boldsymbol{w}^{\prime} \in \mathbb{S}^{n-1}$ with $w_n^{\prime} \neq 0$, we can rewrite $X_n$ with $X_1,...,X_{n-1}$:
\begin{align*}
X_n=\frac{1}{w_n^{\prime}}\left(\boldsymbol{u}-\sum_{i=1}^{n-1} w_i^{\prime} X_i\right),
\end{align*}
where $\boldsymbol{u}:=\boldsymbol{w}^{\prime T} M_n$. Indeed, since $X_n \cdot \boldsymbol{y}$ is irrelevant to how $\boldsymbol{w}^{\prime}$ is chosen,
thus, for the event $\left\{s_n\left(M_n\right) \leq \eta\right\} \wedge\left\{\left\|\boldsymbol{a}^{\prime}{ }_{M_n}\right\|_{\infty}=\left|a_n^{\prime}\right|\right\}$ to occur, we must necessarily have
\begin{align}\label{pro31}
\left|X_n \cdot \boldsymbol{y}\right|  =\inf _{\boldsymbol{w}^{\prime} \in \mathbb{S}^{n-1}, w_n^{\prime} \neq 0} \frac{1}{\left|w_n^{\prime}\right|}\left|\boldsymbol{u} \cdot \boldsymbol{y}-\sum_{i=1}^{n-1} w_i^{\prime} X_i \cdot \boldsymbol{y}\right|,
\end{align}
then, using the inequality Cauchy-Schwarz inequality for the right-hand side of (\ref{pro31}) and setting $\boldsymbol{w}^{\prime}=\boldsymbol{a}^{\prime}{ }_{M_n}$, we have
\begin{align*}
\left|X_n \cdot \boldsymbol{y}\right|&\leq\frac{1}{\left|a_n^{\prime}\right|}\left(\left\|\boldsymbol{a}_{M_n}^{\prime T} M_n\right\|_2\|\boldsymbol{y}\|_2+\left\|\boldsymbol{a}_{M_n}^{\prime}\right\|_2\left(\sum_{i=1}^{n-1}\left(X_i \cdot \boldsymbol{y}\right)^2\right)^{1 / 2}\right)\\
&\leq\eta\sqrt{n}\left(\|\boldsymbol{y}\|_2+\left\|\boldsymbol{a}_{M_n}^{\prime}\right\|_2\right) \leq 2 \eta \sqrt{n},
\end{align*}
where the second inequality is due to the assumption that event $\{s_n\left(M_n\right) \leq \eta\} \wedge\{\left\|\boldsymbol{a}^{\prime}_{M_n}\right\|_{\infty}=\left|a_n^{\prime}\right|\}$ occurs. It follows, by definition, that the probability in (\ref{pro30}) is bounded by $\mathcal{L}\left(X_n \cdot \boldsymbol{y}, 2 \eta \sqrt{n}\right)$, and using Lemma \ref{lcd}, by
\begin{align*}
\mathcal{L}\left(X_n \cdot \boldsymbol{y}, 2 \eta \sqrt{n}\right) \lesssim \eta \sqrt{\frac{n}{\mu_{n}}}+\exp (-2\sqrt{n}\mu_{n}).
\end{align*}
which completes the proof.
\end{proof}

\end{document}
